\begin{abstract}
Medidas de conservação de água podem, involuntariamente, elevar a salinidade do esgoto tratado ao reduzir o volume de água que dilui a mesma carga de sais. Este trabalho investiga a tendência de longo prazo de Sólidos Dissolvidos Totais (TDS) no efluente da Los Angeles--Glendale Water Reclamation Plant (LAGWRP), a partir de 182 meses de dados reais do portal eSMR (2011--2026), com uma bateria de 35 métodos sob um mesmo protocolo de validação temporal honesta. Uma tendência de alta estatisticamente significativa (2,9 a 3,9 mg/L/ano) aparece nos métodos de ajuste linear global, mas não sobrevive a recortes temporais nem a certas correções de autocorrelação -- em vez disso, a série se comporta de forma \textbf{cíclica por regime}, com picos coincidindo com as duas secas recentes da Califórnia (2012--2016 e 2020--2022). Esse reenquadramento foi confirmado com uma variável externa independente (o Índice de Severidade de Seca de Palmer do NOAA) e decomposto via normalização por vazão (WRTDS): a tendência bruta praticamente desaparece após descontar o efeito da vazão (p=0,064). Um balanço de massa não-circular confirma o mecanismo causal -- a carga de sal caiu 3,20\%/ano e a vazão caiu 3,85\%/ano no mesmo período --, ou seja, \textbf{diluição, não aumento de sal}. Como a extrapolação linear se mostrou fisicamente implausível nesta série dirigida por seca, a previsão para +10/+15/+20 anos é apresentada como uma faixa condicionada a cenários climáticos via Monte Carlo (594 a 718 mg/L em +20 anos, seco a úmido), não como um valor único. TDS correlaciona-se positivamente com Amônia mesmo após remoção de tendência comum ($r=0{,}175$, p=0,018), consistente com inibição da nitrificação por salinidade; a correlação com BOD é negativa e não significativa em todos os tratamentos de dado censurado testados -- o oposto do previsto, reportado como tal. Os achados são compatíveis com o mecanismo de conservação descrito na literatura \cite{schwabe2020unintended}, mas um estudo independente com 26 estações da mesma região \cite{scsc2018tds} atribui a esse mecanismo um papel secundário (~12\% da variabilidade) frente à qualidade da água de origem (~88\%), variável ausente do conjunto de dados usado aqui e registrada como limitação central.
\end{abstract}

\noindent \textbf{Palavras-chave:} Salinidade de efluente, séries temporais cíclicas, WRTDS, balanço de massa, projeção por cenários climáticos.
